\section{Design Tradeoffs and Engineering Decisions}

Ghostext is not designed as a maximal-capacity prototype. It is designed as a controlled engineering system where each component has an explicit failure surface and traceable rationale. This section explains those choices and why they were prioritized for the current project phase.

\subsection{Capacity vs. Determinism}

A common optimization target in linguistic steganography is bits per token. In principle, one can push utilization by broadening candidate sets or by applying stronger distribution adjustments. In practice, these moves increase synchronization sensitivity. More candidates mean more opportunities for retokenization mismatch and finite-precision edge cases. More aggressive adjustment can increase divergence from baseline model behavior or create mode-switching behavior across contexts.

Ghostext deliberately places determinism first. The present system treats decode correctness as a hard requirement and accepts lower utilization when needed. This decision is aligned with the current stage objective: establish a trustworthy protocol substrate before adding aggressive optimization layers.

\subsection{Naturalness vs. Invertibility}

There is also tension between linguistic naturalness and strict invertibility. Free-form generation heuristics that improve fluency can break deterministic replay unless the same heuristic is exactly replayable under identical states. Conversely, rigid deterministic rules can slightly reduce surface fluency.

Ghostext chooses a split strategy. The packet-carrying prefix is generated under strict invertible rules. Optional natural tail tokens are added only after payload resolution, so they can improve ending quality without affecting decode correctness. The decoder consumes only the payload-resolving prefix and reports any trailing tokens separately.

\subsection{Tokenizer Fidelity vs. Throughput}

Retokenization-stability checks are computationally nontrivial because they require render-and-retokenize validation for candidate paths. Skipping them would improve throughput but increase risk of silent path divergence.

Ghostext keeps the guard available and enables strict behavior in the candidate policy. In low-risk settings, users may choose looser settings for speed; in higher-risk settings, stability checks remain the safer default.

This is one of several places where Ghostext favors explicit policy knobs over hidden heuristics.

\subsection{Truncation Practicality vs. Distribution Fidelity}

Top-$p$/max-$k$ truncation is operationally attractive because it bounds per-step work and avoids very long probability tails. However, truncation necessarily modifies the effective distribution used for embedding. NLS and OD-Stega both discuss this family of tradeoff~\cite{ziegler2019neural,huang2026odstega}.

Ghostext currently uses deterministic truncation and renormalization with explicit parameters that are part of the synchronization fingerprint. The important design point is not that truncation is universally optimal, but that its exact form is protocol-visible and reproducible.

\subsection{Retry Strategy: Reliability Benefit and Security Cost}

Retries with fresh packet randomness can rescue encoding in low-entropy or unstable-tokenization regions. This increases usability and success rate. At the same time, retries may increase latency variance and can, in principle, alter outward traffic patterns.

Ghostext therefore keeps retry behavior bounded and explicit: maximum attempts are configurable, failures are surfaced, and retry-causing conditions are class-labeled. This makes operational behavior auditable and easier to evaluate in future experiments.

\subsection{Strict Mismatch Rejection vs. Cross-Version Compatibility}

Fingerprint checks reject mismatched model/backend/prompt settings, which strongly reduces silent corruption risk. The downside is reduced tolerance to cross-version deployment drift.

Ghostext currently accepts this downside because ambiguity-tolerant decoding would require significantly more complex compatibility logic and likely introduce security ambiguity. For the current phase, strict rejection is the cleaner and safer policy.

\subsection{Interpretability vs. Auto-Tuning}

Many practical ML systems improve performance through adaptive or learned auto-tuning layers. In steganography, these layers can improve utility but can also blur interpretability and complicate reproducibility.

Ghostext currently favors interpretable deterministic components with explicit parameters. This makes reasoning, testing, and debugging easier during protocol stabilization. Auto-tuning can be explored later once baseline behavior is fully characterized.

\subsection{What These Tradeoffs Mean for ``Almost Perfect''}

Taken together, these decisions define the practical meaning of ``almost perfect'' in Ghostext's current stage: deterministic, auditable, and conservative behavior first; higher-capacity optimization later. The phrase should therefore be read as an engineering trajectory under explicit constraints, not as an assertion that all tradeoffs have already been optimized.
